% Nouveau EGA V (cote 25): archive cote 25, site page 8 (PDF page 9).
% typescript. Draft compared: none.
% AI-transcribed by Claude Opus 5.5 (claude-opus-5-5, adaptive thinking, effort high), Message Batches API, batch ega25b (msgbatch_01A7n5YuXubUPNK1zsepi8RG), 2026-11-13; not yet checked by eye.
\scanpage{25/8}
\margin{\illegible{1}}

\margin{$|$}Un oubli terminologique: lorsque f est une immersion, on appelle habituel%
lement Y\hand{$_\zeta$} une \uline{section hyperplane} de X (relative à l'immersion projective
f et à l'hyperplan H  ); il n'y a pas de raison de ne pas étendre cette
terminologie au cas f quelconque.

20.2. \uline{Etude de la section hyperplane\xout{m} générique}: \hand{\uline{propriétés locales}} .

Rappelons que maintenaht S=Spec(k), k un corps. Si \hand{$\eta$} est un point
de $\overset{\hand{$\vee$}}{P}$ , \xout{xxxx} et si \hand{$\zeta$}: Spec(\hand{$\eta$})$\to\overset{\hand{$\vee$}}{P}$ est le morphisme canonique, nous poserons
aussi H\hand{$_\eta$} , Y\hand{$_\eta$} , G\hand{$_\eta$} au lieu de H\hand{$_\zeta$} , Y\hand{$_\zeta$} , G\hand{$_\zeta$} . Dans le présent nº, \hand{$\eta$}
désignera toujours le point générique de $\overset{\hand{$\vee$}}{P}$ .
\uline{Proposition} 2.1. Supposons X irréductible\inserted{\deleted{\illegible{1}}}. Alors Y\hand{$_\eta$} est irr. \inserted{(\deleted{\illegible{1}})} ou vide. ,
et dans le premier cas il domine X. \inserted{En tous cas $\unsure{\overline{Y}}$ est \unsure{$_\eta$} irréductible.}
En effet, comme \uline{H}$\longrightarrow$ P est un fibré projectif, il en est de même de
\uline{Y}$\longrightarrow$ X, ce qui implique que \uline{Y} est irréductible X l'\unsure{étant}. Danc la fibre
générique Y   de \uline{Y} sur $\overset{\hand{$\vee$}}{P}$ est irréductible ou vide, et dans le premier cas,
son point générique est point générique de \uline{Y} , donc au dessus de celui de
X, cqfd.
\uline{Proposition 2}.2. Soit Z une partie de P. Alors son image inverse Z\hand{$_\eta$} dans
H\hand{$_\eta$} est vide si et seulement si tout point de Z est fermé. En particulier,
si Z est constructible, alors Z\hand{$_\eta$} = $\emptyset$ sss Z est fini.

On peut supposer Z réduit à un seul point z , et on est ramené à
prouver que l'\unsure{imalge|image} de H dans P est formé \inserted[typed]{exactement} des points non fermés de P.
Désignant par X l'adhérence de z, et utilisant 2.1. , on est ramené à
prouver que \over{X}{\unsure{Y}}\hand{$_\eta$} = $\emptyset$ sss X est \xout{\illegible{3}} \inserted[typed]{fini} (X étant un sous-schém%
fermé \xout{xxxx} de P). Remplaçant X par X$_{\text{k(\hand{$\eta$})}}$\hand{$\hookrightarrow$}P$_{\text{k(\hand{$\eta$})}}$ , \xout{xxxx} \inserted[typed]{le ``il faut''} résulte du fait
\margin{$|$}suivant, pour lequel nous devons avoir une référence, et qui mérite d'être
réex\unsure{mplicité|plicité} ici en lemme: si Y est une sectionhyperplane \uline{quelconque} de
X et si Y\hand{\unsure{$_\eta$}}=$\emptyset$ , alors X est fini ( en effet, X $\subset$ P - H est affine, et
projectif ...). Le "il suffit" est immédiat, par exemple en notant que
\uline{Y} est un fibré projectif de dimension relative n-1 sur X (\xout{x}n étant la
dim relative de P \xout{xxxxS)} et $\overset{\hand{$\vee$}}{P}$ sur S), donc X étant fini sur k, \uline{Y} est de

\mspage{\hand{8}}
